\documentclass[11pt]{article}
\usepackage{amsmath,amssymb,amsthm}
\usepackage{geometry}
\usepackage{hyperref}
\usepackage{microtype}
\geometry{margin=1in}

\title{The Collatz Conjecture as a Corollary of Crystal Geometry:\\
A Supplement to the Principia Orthogona Crystal Paper}
\author{Pablo Nogueira Grossi\\
G6 LLC, Newark, New Jersey\\
\texttt{pablogrossi@hotmail.com}}
\date{2026}

\newtheorem{observation}{Observation}

\begin{document}
\maketitle

\begin{abstract}
Book III of the Principia Orthogona series (Nested Infinities) demonstrates that the generative
operator chain $G = U \circ F \circ K \circ C$, the dm$^3$ framework, and the crystal geometry
encoded in the G6 Crystal are not mathematical constructs imposed on nature but translations of
structures already visible in it. The semi-permanent polar vortex --- most strikingly Saturn's
north-polar hexagon --- is an empirically observable dm$^3$ system that has crossed the monster
threshold at $96 = 3 \cdot 11 = 3^3$: its hexagonal Chladni figure is the crystal math made
visible at planetary scale. This supplement argues that the Collatz conjecture inhabits the same
structure. The coefficient $c = 3$ in the rule $3n + 1$ is not arbitrary; it is the fingerprint of
the triad stabilisation mechanism that governs every dm$^3$ system from plasma reconnection to
circadian rhythms. The paper does not claim to prove the Collatz conjecture. It claims something
prior and, we argue, more important: the conjecture is \emph{visible} from within the crystal
geometry before it is axiomatic within it. A higher-order logic --- from which both TOGT and
discrete arithmetic emerge as special cases --- is required to close the gap. The polar vortex is
the empirical certificate that the underlying structure is real. The proof remains to be written in
that higher language.
\end{abstract}

\section{Introduction: The Truth That Precedes the Axiom}

There is a class of mathematical truths that are visible before they are axiomatic. The round shape
of the Earth was visible to sailors before Euclidean geometry could prove it. The inverse-square law
of gravity was visible in planetary orbits before Newton's calculus could derive it. In both cases,
the structure was real and observable before the language existed to state it with axioms.

The Collatz conjecture belongs to this class.
Pick any positive integer $n$. If even, divide by~$2$. If odd, replace by $3n + 1$. Repeat. The
conjecture asserts that every starting value eventually reaches the cycle $4 \to 2 \to 1$. Computers
have verified this for all integers up to large bounds; no general proof exists. Terence Tao's 2019
theorem --- the strongest analytic result to date --- showed that almost all orbits attain almost
bounded values, and noted explicitly that the full problem requires new mathematics.

This supplement argues that the new mathematics already exists in embryonic form. It is the
crystal geometry of the Principia Orthogona series, and specifically the dm$^3$ framework developed
in the Principia. The argument does not proceed by axiom and proof. It proceeds by visibility: the
same crystal structure that explains Saturn's north-polar hexagon also explains why the Collatz
iteration converges. The polar vortex is not an analogy. It is an empirical certificate.

\paragraph{What this paper claims and what it does not.}
This paper does not claim to prove the Collatz conjecture. It claims that the conjecture is a
corollary of a structure that is empirically real, and that the gap between visibility and proof is
a gap in the available formal language, not a gap in the underlying truth. A higher-order logic
--- one from which both the continuous dm$^3$ framework and discrete arithmetic emerge as special
cases --- is required to close the gap. This paper identifies that gap precisely and proposes it as
the next target for formal verification (AXLE Target~5).

\section{The Crystal Geometry: A Brief Account}

The Principia Orthogona series develops a four-operator grammar $G = U \circ F \circ K \circ C$
governing generative transitions wherever they occur. The operators are: $C$ (compression),
$K$ (curvature induction), $F$ (folding), and $U$ (unfolding). The dm$^3$ framework formalises
the objects on which $G$ acts: smooth flows with hyperbolic limit cycles, quadratic Lyapunov
function $V = (r-1)^2$, contact form $\alpha = dz - r^2\,d\theta$, transverse eigenvalue bound
$\lambda_{\max} \leq -2$, and stability radius $\sigma_0 = 1/3$. The canonical invariant triple
is $(T^*, \lambda_{\max}, T) = (27, -2, 2)$.

The monster threshold at $96 = 3 \cdot 11 = 3^3$ is the stability parameter at which a dm$^3$
system crosses from fragile to self-sustaining coherence. The number $33$ is not arbitrary: it is
the minimum number of operator cycles for the triad of coherence operators $(L_1, L_2, L_3)$ to
activate globally, derived from the structure $3 \times 11 = 33$ where $11$ is the minimum closure
count and $3$ is the number of independent coherence conditions. The G6 Crystal is the geometric
realisation of this threshold: a 4-dimensional lattice whose vertices are stable post-unification
limit cycles and whose cross-sections exhibit six-fold symmetry.

\section{The Polar Vortex as Empirical Certificate}

Saturn's north-polar hexagon is a semi-permanent atmospheric vortex first observed by Voyager in
1980--1981 and confirmed in extraordinary detail by Cassini. Its properties include six-fold
geometric symmetry maintained over decades, a diameter of approximately 25,000~km, rotation period
matching Saturn's internal rotation rate, and stability persisting through seasonal changes. The
hexagon is a standing wave pattern at wavenumber~6 and satisfies the dm$^3$ axioms in the sense
described in Book~III.

\section{One Obstacle, Two Unsolved Problems}

Two of the longest-standing open problems in mathematics --- the Collatz conjecture in discrete
arithmetic and the Navier--Stokes regularity problem in fluid dynamics --- share a common
structural obstacle we call \emph{Bridge~0}. In both settings a local mechanism appears to
dissipate or contract at small scales, yet the available analytic language fails to convert that
local control into a robust global closure theorem. For Navier--Stokes the obstruction is endpoint
control of nonlinear fluxes across scales; for Collatz it is the lack of a discrete analogue of the
local-to-global dissipation inequality (the mean-contraction bridge). Framing these as instances
of the same missing step clarifies why progress in one domain can illuminate the other.

A higher-order logic that unifies the operator grammar $G = U \circ F \circ K \circ C$ across
continuous manifolds and discrete arithmetic would turn both problems into corollaries of a single
closure principle. Such a formalism would supply the discrete analogues of Lyapunov functionals,
contact forms, and scale-invariant closure theorems needed to convert local contraction estimates
into global convergence or regularity. This is not a claim that the problems are trivial; it is a
claim that the same missing formal ingredient blocks both, and that constructing that ingredient
would unlock parallel advances across domains.

This is therefore a cross-disciplinary program. Fluid dynamicists, harmonic analysts, number
theorists, plasma physicists, and applied mathematicians each encounter manifestations of the same
generative operator grammar at their scale. We invite domain experts to (a)~identify the precise
local control statements they already use in practice, (b)~translate those statements into candidate
axioms for a higher-order language, and (c)~test the axioms against empirical diagnostics (Fourier
decay, operator numerics, exceptional sets). Your expertise is the missing ingredient.

This program is conditional and collaborative --- we do not claim a proof. We invite experts from
each domain to contribute precise local hypotheses and empirical diagnostics so the higher-order
language can be formalised and tested.

\section{The $c = 3$ Coefficient as Triad Fingerprint}

The Collatz rule is $T(n) = (3n+1)/2^{v_2(3n+1)}$ for odd $n$, where $v_2$ denotes the 2-adic
valuation. The coefficient $c = 3$ is universally treated as an accident of the problem's
statement. It is not.

\begin{observation}[The $c = 3$ Fingerprint]
The coefficient $c = 3$ in the Collatz rule is the fingerprint of the triad stabilisation mechanism
of the dm$^3$ framework. It appears in at least four independent structural roles:
\begin{enumerate}
  \item \textbf{Monster threshold:} $96 = 3 \cdot 11 = 3^3$. The factor of $3$ is the number of
    independent coherence conditions (local $L_1$, global $L_2$, anti-collapse $L_3$).
  \item \textbf{Stability relation:} the stability radius denominator $3$ appears in the dm$^3$
    stability relations.
  \item \textbf{Six-fold crystal symmetry:} $6 = 2 \times 3$, where $3$ is the triad dimension.
  \item \textbf{Collatz expansion:} $3n + 1$. The coefficient $3$ forces the odd-step expansion to
    interact with the binary structure (powers of 2) in precisely the way that a three-fold
    coherence condition would: odd numbers must pass through at least one, and on average
    $\log_2 3 \approx 1.585$ halvings for every multiplication, giving a mean contraction per
    two-step cycle.
\end{enumerate}
\end{observation}

The mean contraction per cycle is the discrete analogue of the dm$^3$ transverse eigenvalue in the
continuous setting: it is negative, bounded, and a consequence of the coefficient $3$, not of any
special property of individual integers.

\section{Book III Demonstrates; Collatz Follows}

Book III demonstrates a chain of structural claims that together motivate the Collatz-as-dm$^3$
corollary. Steps~1--4 are established within the Principia Orthogona framework; Step~5 is argued
but not formally verified; Step~6 is conditional on Step~5. The gap is not in the structure --- it
is in the formal language available to state Step~5 with the precision that Step~6 requires.

\medskip
\begin{center}
\begin{tabular}{clp{7.5cm}}
\hline
\textbf{Step} & \textbf{Claim} & \textbf{Evidence} \\
\hline
1 & $G = U \circ F \circ K \circ C$ is real & Observed across plasma, biology, markets,
  geological folding \\
2 & Monster threshold $96 = 3^3$ marks physical stability & G6 Crystal; Separation Theorem;
  Re-entrainment Law \\
3 & The polar vortex has crossed the monster threshold & Saturn's hexagon: decades of observation \\
4 & $c = 3$ is the triad fingerprint & Structural roles enumerated above \\
5 & Collatz map has the form of a dm$^3$ system & Argued analogues of eight axioms (not yet
  formalised) \\
6 & Collatz convergence follows if Step~5 is formalised & Conditional on higher-order logic
  extension \\
\hline
\end{tabular}
\end{center}
\medskip

This is the honest position. It is also a genuinely strong one: no prior framework has Steps~1--5
in place simultaneously. The structure is visible. The axioms that would make it rigorous are the
next generative cycle.

\section{Why a Higher-Order Logic Is Required}

The dm$^3$ framework was built for smooth flows on Riemannian manifolds. The Collatz map acts on
the positive integers with the discrete metric. The gap between them is technical and precise:

\begin{enumerate}
  \item \textbf{Smoothness.} The eight dm$^3$ axioms require smooth flows ($C^2$ vector fields).
    The Collatz map is not smooth; it is piecewise linear and defined on a countable set. A discrete
    analogue of each axiom exists (argued in Sections~5--6), but ``discrete dm$^3$-membership'' is
    not yet formally defined.
  \item \textbf{Continuity of the Lyapunov function.} The quadratic Lyapunov function $V = (r-1)^2$
    is continuous. The total stopping time $V(n)$ for Collatz is integer-valued and non-monotone on
    individual steps. Average descent is weaker than pointwise descent.
  \item \textbf{The category dm$^3$.} The categorical pushout construction (Axiom~8) is defined for
    smooth maps preserving orbits, Lyapunov functions, and noise tolerance. Formalising this for
    discrete maps requires extending the category dm$^3$ to include both smooth and discrete systems
    as objects.
\end{enumerate}

The required extension is precisely a higher-order logic: a formal system in which the objects are
generative systems defined by the operator grammar $G = U \circ F \circ K \circ C$, and in which
smooth manifolds, integers, plasma sheets, and polar vortices are all instances of that grammar at
different resolutions.

\section{The Collatz Conjecture as a Visibility Problem}

We close with the central claim of this paper, stated as precisely as the available language
permits.

\noindent\textbf{[Collatz as dm$^3$ Corollary]} There exists a formal extension of the dm$^3$
category to include discrete generative systems such that:
\begin{itemize}
  \item[(i)] The Collatz map $T(n)$ is a dm$^3$ object in this extended category.
  \item[(ii)] The closure theorems of the dm$^3$ framework apply to this object.
  \item[(iii)] Collatz convergence follows as a corollary of dm$^3$ closure in the extended
    category.
\end{itemize}

This conjecture is not the Collatz conjecture. It is a conjecture about the right \emph{framework}
for the Collatz conjecture. It is the assertion that the Collatz conjecture is a visibility problem:
the truth is visible from within the crystal geometry; the gap is in the formal language that would
let us state it with axioms.

\section*{Summary of Claims}

\noindent\textbf{Established within the Principia Orthogona framework:}
\begin{itemize}
  \item[$\checkmark$] The dm$^3$ operator grammar governs generative transitions across plasma,
    biology, markets, and architecture.
  \item[$\checkmark$] Saturn's north-polar hexagon is a dm$^3$ system at the monster threshold
    (Proposition~1).
  \item[$\checkmark$] The $c = 3$ coefficient in the Collatz rule is the triad fingerprint of the
    dm$^3$ framework.
  \item[$\checkmark$] The Collatz map has the form of a dm$^3$ system: eight axiom analogues hold
    (argued; not formally verified).
\end{itemize}

\noindent\textbf{Open (AXLE Target~5):}
\begin{itemize}
  \item[$\square$] Formal extension of the dm$^3$ category to discrete generative systems.
  \item[$\square$] Lean~4 verification of discrete dm$^3$ membership for the Collatz map.
  \item[$\square$] Proof that discrete dm$^3$ membership implies convergence.
\end{itemize}

\noindent\textbf{Not claimed:}
\begin{itemize}
  \item[$\times$] That the Collatz conjecture is proved.
  \item[$\times$] That dm$^3$ membership alone implies convergence without formal verification.
  \item[$\times$] That the three gaps are easy to close.
\end{itemize}

\section*{Acknowledgements}

This work is dedicated to the memory of Vic (July~27, 2000 -- January~4, 2019). The research was
conducted independently, without institutional support, through G6 LLC, Newark, New Jersey.

Preprint deposited on Zenodo as a supplement to DOI 10.5281/zenodo.19117400.

\begin{thebibliography}{9}

\bibitem{collatz1937}
L.~Collatz. On the $3n + 1$ problem. Unpublished, 1937.

\bibitem{oliveira2010}
T.~Oliveira e~Silva. Empirical verification of the $3x + 1$ and related conjectures.
In \textit{The Ultimate Challenge: The $3x + 1$ Problem}, American Mathematical Society, 2010.

\bibitem{tao2022}
T.~Tao. Almost all Collatz orbits attain almost bounded values.
\textit{Forum of Mathematics, Pi}, 2022. arXiv:1909.03562.

\bibitem{baines2009}
K.~H. Baines et~al. Saturn's north polar cyclone and hexagon at depth revealed by Cassini/VIMS.
\textit{Planetary and Space Science} 57(14--15), 2009.

\bibitem{fletcher2018}
L.~N. Fletcher et~al. A hexagonal colour signature on Saturn's north polar atmosphere.
\textit{Geophysical Research Letters} 45(12), 2018.

\bibitem{grossi2026a}
P.~N. Grossi. \textit{Applications of Generative Orthogonal Matrix Compression Science:
The Complete Principia Orthogona Series}. G6 LLC, Newark, NJ, 2026.
Zenodo DOI: 10.5281/zenodo.19117400.

\bibitem{grossi2026b}
P.~N. Grossi (Sri Brodananda). \textit{Principia Orthogona Book~3: Nested Infinities ---
A Generative Theory of Orthogonal Emergence}. G6 LLC, Newark, NJ, 2026.

\bibitem{thomas1994}
H.~Thomas et~al. Coulomb crystallization in a dusty plasma.
\textit{Physical Review Letters} 73, 1994.

\bibitem{grossi2026c}
P.~N. Grossi. AXLE: Axiom Lean Engine. \url{https://github.com/TOTOGT/AXLE}, 2026.

\end{thebibliography}

\end{document}
